\documentclass[11pt]{article}
\usepackage{amsmath}
\usepackage{amssymb}
\usepackage{geometry}
\usepackage{bm}
\usepackage{graphicx}
\usepackage{hyperref}
\usepackage{float}
\usepackage{booktabs}
\geometry{margin=1in}

\newcommand{\vp}{\varphi}
\newcommand{\bS}{\mathbf{S}}
\newcommand{\bxi}{\boldsymbol{\xi}}

\title{Geometric Basin Connectivity in the Large-$b$ Regime:\\
  Exact Spherical Analysis}
\author{}
\date{}

\begin{document}
\maketitle

\begin{abstract}
We study geometric connectivity of pattern basins in the
LSR associative memory when the nonlinearity parameter scales as
$b = \beta N$ with $\beta > 0$ fixed.
The analysis proceeds in two stages: first purely geometric
(do the spherical caps overlap at all?), then energetic
(can the state traverse the corridor?).
For $b = \beta N$, we show that energy barriers between
geometrically overlapping basins are $O(1/\beta)$,
so geometry is the sole bottleneck.
We then compute the expected number of geometric neighbours
$\lambda$ using the \emph{exact} spherical overlap distribution
and compare with the Gaussian (Mill's ratio) approximation.
The two predictions differ by a super-exponential factor:
the Gaussian gives a percolation threshold
$\alpha_c^{\rm Gauss} \to 1/2$,
while the exact sphere gives
$\alpha_c^{\rm sphere} \approx \tfrac{1}{2}\ln(\beta N/8)$,
which diverges with~$N$.
We conclude that for $b = \beta N$ and large $N$,
every basin is geometrically isolated at any fixed load~$\alpha$.
The phase transition at $\alpha = 1/2$ in the exponential
associative memory is driven by signal-to-noise competition
in the softmax energy, not by basin connectivity.
\end{abstract}


%% ==================================================================
\section{The model}
\label{sec:model}

The Log-Sum-ReLU (LSR) associative memory stores $M$ patterns
$\bxi^1, \ldots, \bxi^M$ drawn uniformly on the sphere
$S^{N-1}(\sqrt{N})$ (i.e.\ $|\bxi^\mu|^2 = N$).
The energy of a state $\bS \in S^{N-1}(\sqrt{N})$ is
\begin{equation}\label{eq:E}
  E(\bS) = -\frac{N}{b}\sum_{\mu=1}^{M}
    \ln\bigl[\,1 - b + b\,\vp_\mu(\bS)\,\bigr]_+\,,
\end{equation}
where $\vp_\mu(\bS) = (\bxi^\mu \cdot \bS)/N$ is the overlap
between the state and pattern~$\mu$,
$[\cdot]_+ = \max(0, \cdot)$,
and $b > 0$ is the nonlinearity parameter.
The number of patterns scales exponentially with dimension:
$M = e^{N\alpha}$, where $\alpha > 0$ is the load.

Throughout this note we work in the regime
\begin{equation}\label{eq:b_regime}
  b = \beta N, \qquad \beta > 0 \;\text{fixed},
  \qquad \beta N \gg 1.
\end{equation}


\subsection{Hard-wall threshold and basins}

The argument of the logarithm in~\eqref{eq:E} is
$1 - b + b\,\vp_\mu = b(\vp_\mu - \vp_c)$,
where
\begin{equation}\label{eq:phic}
  \vp_c \;=\; \frac{b-1}{b} \;=\; 1 - \frac{1}{b}
  \;=\; 1 - \frac{1}{\beta N}\,.
\end{equation}
Pattern $\mu$ contributes to the energy only when
$\vp_\mu(\bS) > \vp_c$; otherwise $[\cdots]_+ = 0$.
If \emph{no} pattern is active, the sum is empty and
$E = +\infty$ (the state is energetically forbidden).

The \textbf{basin} of pattern $\mu$ is the spherical cap
\begin{equation}\label{eq:basin}
  X_\mu = \bigl\{\, \bS \in S^{N-1}(\sqrt{N})
  \;:\; \vp_\mu(\bS) > \vp_c \,\bigr\}.
\end{equation}
The boundary $\vp_\mu = \vp_c$ is an infinite energy barrier
($\ln 0^+ = -\infty$ in the energy, so $E \to +\infty$).
Within a single basin, the state is permanently confined.


\subsection{Angular radius of a basin}

The condition $\vp_\mu > \vp_c$ is equivalent to
$\cos\theta < \vp_c$, where $\theta$ is the angle between
$\bS$ and $\bxi^\mu$.  The cap has angular radius
\begin{equation}\label{eq:thetac}
  \theta_c = \arccos\vp_c
  = \arccos\!\Bigl(1 - \frac{1}{\beta N}\Bigr)
  \approx \sqrt{\frac{2}{\beta N}}
  \qquad (\beta N \gg 1),
\end{equation}
using $\arccos(1 - x) \approx \sqrt{2x}$ for $x \ll 1$.
So each basin is a tiny cap whose angular radius shrinks as
$1/\sqrt{\beta N}$.


%% ==================================================================
\section{Geometric overlap of two basins}
\label{sec:geom_overlap}

\subsection{Criterion for cap intersection}

Two spherical caps of equal angular radius $\theta_c$ overlap
if and only if the angle $\theta_{\mu\nu}$ between their
centres satisfies $\theta_{\mu\nu} < 2\theta_c$.
In terms of the mutual overlap
$q = (\bxi^\mu \cdot \bxi^\nu)/N = \cos\theta_{\mu\nu}$,
the caps overlap when
\begin{equation}\label{eq:qgeom}
  q > q_{\rm geom}
  \;\equiv\; \cos(2\theta_c)
  = 2\cos^2\theta_c - 1
  = 2\vp_c^2 - 1.
\end{equation}
Substituting $\vp_c = 1 - 1/(\beta N)$:
\begin{equation}\label{eq:qgeom_explicit}
  q_{\rm geom}
  = 2\Bigl(1 - \frac{1}{\beta N}\Bigr)^{\!2} - 1
  = 1 - \frac{4}{\beta N} + \frac{2}{\beta^2 N^2}\,.
\end{equation}
This is extremely close to~$1$: the two patterns must be
nearly identical for their basins to touch.

If $q < q_{\rm geom}$, the caps are disjoint and no continuous
path connects them without leaving all basins (hitting $E = +\infty$).
No amount of thermal energy can bridge a gap where no corridor
exists.


\subsection{The corridor between overlapping basins}

When $q > q_{\rm geom}$, the caps $X_\mu$ and $X_\nu$ share
a non-empty intersection: the \textbf{corridor}
$X_\mu \cap X_\nu$.  Along any path through the corridor,
at least two patterns are active, and the energy remains finite.

To compute the energy along the corridor, consider the geodesic
(great-circle arc) from $\bxi^\mu$ to $\bxi^\nu$ on
$S^{N-1}(\sqrt{N})$.
Parametrise it by the angle $\psi \in [0, \theta_{\mu\nu}]$
from $\bxi^\mu$:
\begin{equation}\label{eq:geodesic}
  \bS(\psi)
  = \frac{\sin(\theta_{\mu\nu} - \psi)}{\sin\theta_{\mu\nu}}\,\bxi^\mu
  + \frac{\sin\psi}{\sin\theta_{\mu\nu}}\,\bxi^\nu.
\end{equation}
The overlaps along this geodesic are
\begin{align}
  \vp_\mu(\psi) &= \frac{\bxi^\mu \cdot \bS(\psi)}{N}
  = \frac{\sin(\theta_{\mu\nu} - \psi)}{\sin\theta_{\mu\nu}}
  + q\,\frac{\sin\psi}{\sin\theta_{\mu\nu}}\,,
  \label{eq:phi_mu_psi} \\[4pt]
  \vp_\nu(\psi) &= \frac{\bxi^\nu \cdot \bS(\psi)}{N}
  = q\,\frac{\sin(\theta_{\mu\nu} - \psi)}{\sin\theta_{\mu\nu}}
  + \frac{\sin\psi}{\sin\theta_{\mu\nu}}\,.
  \label{eq:phi_nu_psi}
\end{align}
At $\psi = 0$ (centre of $\mu$): $\vp_\mu = 1$, $\vp_\nu = q$.
At $\psi = \theta_{\mu\nu}$ (centre of $\nu$):
$\vp_\mu = q$, $\vp_\nu = 1$.


%% ==================================================================
\section{Energy barrier along the geodesic}
\label{sec:barrier}

\subsection{Energy of a single active pattern}

When only pattern $\mu$ is active (with overlap $\vp_\mu > \vp_c$),
its contribution to the energy is
\begin{equation}\label{eq:E_single}
  E_\mu = -\frac{N}{b}\,\ln\bigl[1 - b + b\,\vp_\mu\bigr]
  = -\frac{N}{b}\,\ln\bigl[b(\vp_\mu - \vp_c)\bigr].
\end{equation}
At the basin centre ($\vp_\mu = 1$):
\begin{equation}
  E_\mu^{\rm centre}
  = -\frac{N}{b}\,\ln\bigl[b(1 - \vp_c)\bigr]
  = -\frac{N}{b}\,\ln 1
  = 0\,,
\end{equation}
using $b(1 - \vp_c) = b \cdot (1/b) = 1$.

At the basin boundary ($\vp_\mu \to \vp_c^+$):
$E_\mu \to -\frac{N}{b}\ln 0^+ = +\infty$.


\subsection{Energy at the bottleneck}

The energy is highest at the ``bottleneck''---the point on the
geodesic where one pattern has just activated.
By symmetry of the barrier calculation, consider the point $\psi^*$
where pattern $\nu$ just activates: $\vp_\nu(\psi^*) = \vp_c$.
At this point, $\nu$ contributes $\ln(0^+ + 1) = 0$ to the
energy, and the entire energy comes from pattern $\mu$:
\begin{equation}
  E(\psi^*) = -\frac{N}{b}\,\ln\bigl[b(\vp_\mu^* - \vp_c)\bigr],
\end{equation}
where $\vp_\mu^*$ is the overlap with $\mu$ when $\vp_\nu = \vp_c$.

To find $\vp_\mu^*$, we use a coordinate decomposition.
Write $\bxi^\nu = q\,\bxi^\mu + \sqrt{1-q^2}\;\hat{\bxi}^\perp$,
where $\hat{\bxi}^\perp$ is a unit vector (rescaled to
$|\hat\bxi^\perp|^2 = N$) orthogonal to $\bxi^\mu$.
Then at any state $\bS$:
\begin{equation}\label{eq:overlap_decomp}
  \vp_\nu = q\,\vp_\mu + \sqrt{1-q^2}\;\vp_\perp,
\end{equation}
where $\vp_\perp = (\hat{\bxi}^\perp \cdot \bS)/N$.
Setting $\vp_\nu = \vp_c$ and solving for $\vp_\mu$
on the geodesic (where $\vp_\mu^2 + \vp_\perp^2
\leq 1$ with appropriate normalisation),
one finds the overlap with $\mu$ at the $\nu$-activation
boundary.  On the geodesic in the $(\vp_\mu, \vp_\perp)$ plane,
the state satisfies $\vp_\mu = \cos\psi$,
$\vp_\perp = \sin\psi \cdot \sqrt{1 - q^2}/\sin\theta_{\mu\nu}$
(to leading order for small angles).

The result is:
\begin{equation}\label{eq:phi_star}
  \boxed{%
  \vp_\mu^* = q\,\vp_c + \sqrt{(1-q^2)(1-\vp_c^2)}\,.}
\end{equation}
This can be verified geometrically: the point where
the geodesic from $\bxi^\mu$ to $\bxi^\nu$ crosses the
boundary of cap $X_\nu$ (the circle $\vp_\nu = \vp_c$)
has overlap $\vp_\mu^*$ with pattern $\mu$.

\paragraph{Check at $q = q_{\rm geom}$.}
When the caps just touch ($q = 2\vp_c^2 - 1$):
\begin{align}
  \vp_\mu^*
  &= (2\vp_c^2 - 1)\vp_c
    + \sqrt{(1 - (2\vp_c^2-1)^2)(1-\vp_c^2)}
  \nonumber\\
  &= (2\vp_c^2 - 1)\vp_c
    + \sqrt{4\vp_c^2(1-\vp_c^2) \cdot (1-\vp_c^2)}
  \nonumber\\
  &= (2\vp_c^2 - 1)\vp_c + 2\vp_c(1-\vp_c^2)
  \nonumber\\
  &= \vp_c\bigl[2\vp_c^2 - 1 + 2 - 2\vp_c^2\bigr]
  = \vp_c\,.
\end{align}
So $\vp_\mu^* = \vp_c$: both patterns are marginal, confirming
the caps just touch.


\subsection{The barrier height}

The \textbf{barrier} is the energy difference between the
bottleneck and the basin centre:
\begin{equation}\label{eq:barrier}
  \boxed{%
  \Delta E = E(\psi^*) - E_\mu^{\rm centre}
  = -\frac{N}{b}\ln\bigl[b(\vp_\mu^* - \vp_c)\bigr]
  = \frac{N}{b}\,\ln\frac{1}{b\,(\vp_\mu^* - \vp_c)}\,.}
\end{equation}
The barrier diverges as $q \to q_{\rm geom}^+$
(where $\vp_\mu^* \to \vp_c$), confirming that caps which
just barely touch have an infinite barrier.


\subsection{Barrier scaling for $b = \beta N$}

For $b = \beta N$, the prefactor is $N/b = 1/\beta$.
The overlap excess is
\begin{equation}\label{eq:phi_excess}
  \vp_\mu^* - \vp_c
  = q\,\vp_c + \sqrt{(1-q^2)(1-\vp_c^2)} - \vp_c
  = \vp_c(q - 1) + \sqrt{(1-q^2)(1-\vp_c^2)}\,.
\end{equation}
For $q$ slightly above $q_{\rm geom}$,
write $q = q_{\rm geom} + \epsilon$
with $\epsilon \ll 1$.
Using $1 - \vp_c^2 \approx 2/(\beta N)$ and
$1 - q_{\rm geom} \approx 4/(\beta N)$:
\begin{equation}
  \vp_\mu^* - \vp_c
  \;\approx\; \frac{\epsilon}{2\vp_c}
  \;\approx\; \frac{\epsilon}{2}\,.
\end{equation}
The argument of the logarithm in~\eqref{eq:barrier}:
\begin{equation}
  b(\vp_\mu^* - \vp_c)
  \approx \beta N \cdot \frac{\epsilon}{2}
  = \frac{\beta N\,\epsilon}{2}\,.
\end{equation}
For a ``typical'' geometrically overlapping pair,
$\epsilon = O(1/(\beta N))$ (the overlap just barely exceeds
the threshold), so $b(\vp_\mu^* - \vp_c) = O(1)$ and
\begin{equation}\label{eq:barrier_bN}
  \boxed{%
  \Delta E = \frac{1}{\beta}\,
  \ln\frac{1}{\beta N\,(\vp_\mu^* - \vp_c)}
  = O(1/\beta)\,.}
\end{equation}
This is a finite, $N$-independent barrier.
Any temperature $T > 0$ (even $T \ll 1$) overcomes it
on an $O(1)$ timescale, via Kramers-type thermal activation
$\tau \sim e^{\Delta E / T}$ with $\Delta E/T = O(1)$.

\paragraph{Contrast with $b = O(1)$.}
For $b = O(1)$, the prefactor $N/b = O(N)$, giving
$\Delta E = O(N)$: an extensive barrier that grows with
system size and is thermally impassable at any fixed $T$.

\paragraph{Conclusion.}
For $b = \beta N$, \textbf{geometry is the sole bottleneck}.
If two caps overlap ($q > q_{\rm geom}$), the barrier is $O(1/\beta)$
and thermally irrelevant.
If the caps do not overlap ($q < q_{\rm geom}$), no corridor
exists and no amount of thermal energy helps.
We therefore focus entirely on the geometric question:
how many caps overlap?


%% ==================================================================
\section{Overlap distribution on the sphere}
\label{sec:overlap}

For two independent random points on $S^{N-1}(\sqrt{N})$,
the normalised overlap $q = (\bxi^\mu \cdot \bxi^\nu)/N$
is the cosine of the angle between them.
It has the exact marginal density
\begin{equation}\label{eq:fq}
  f(q) = \frac{(1 - q^2)^{(N-3)/2}}{B\!\bigl(\tfrac{1}{2},\,
    \tfrac{N-1}{2}\bigr)}\,,
  \qquad q \in (-1, 1),
\end{equation}
where $B(a,b) = \Gamma(a)\Gamma(b)/\Gamma(a+b)$ is the
beta function.

\paragraph{Derivation.}
Fix $\bxi^\mu$ (say, along the first coordinate axis).
The overlap $q = \xi^\nu_1/\sqrt{N}$ (the first component of
$\bxi^\nu$, rescaled).
For a uniform point on $S^{N-1}(\sqrt{N})$, the projection
onto any axis has the distribution of $\sqrt{N}\cos\Theta$
where $\Theta$ is the polar angle.
The surface area element at polar angle $\Theta$ on
$S^{N-1}$ is proportional to $\sin^{N-2}\Theta\,d\Theta$.
Changing variables to $q = \cos\Theta$:
$dq = -\sin\Theta\,d\Theta$ and
$\sin^{N-2}\Theta = (1-q^2)^{(N-2)/2}$.
Thus $f(q) \propto (1-q^2)^{(N-2)/2}/\sin\Theta
= (1-q^2)^{(N-2)/2}/(1-q^2)^{1/2}
= (1-q^2)^{(N-3)/2}$.
The normalisation constant is $1/B(1/2, (N{-}1)/2)$
by the beta-function integral
$\int_{-1}^{1}(1-q^2)^{(N-3)/2}\,dq = B(1/2, (N{-}1)/2)$.

\paragraph{Moments.}
$\mathbb{E}[q] = 0$ (by symmetry),
$\mathbb{E}[q^2] = 1/(N{-}1) \approx 1/N$.

\paragraph{Gaussian approximation.}
Matching the first two moments, one approximates
$q \sim \mathcal{N}(0, 1/N)$.  This is accurate for
$q = O(1/\sqrt{N})$ but breaks down catastrophically
in the tail near $q = 1$, as we now show.


%% ==================================================================
\section{Tail probability: sphere vs.\ Gaussian}
\label{sec:tail}

The expected number of geometric neighbours depends on
the tail probability $P(q > q_{\rm geom})$.

\subsection{Exact spherical tail near $q = 1$}

For $q_0$ close to $1$, set $q = 1 - t$ with
$\delta = 1 - q_0 \ll 1$:
\begin{equation}
  1 - q^2 = (1-q)(1+q) = t(2-t) \approx 2t
  \qquad (t \ll 1).
\end{equation}
Then:
\begin{align}
  P_{\rm sphere}(q > q_0)
  &= \int_{q_0}^{1} f(q)\,dq
  = C_N \int_0^{\delta}(2t)^{(N-3)/2}\,dt
  \nonumber\\[3pt]
  &= C_N \cdot 2^{(N-3)/2}
    \cdot \frac{\delta^{(N-1)/2}}{(N{-}1)/2}\,,
  \label{eq:Psphere}
\end{align}
where $C_N = 1/B(1/2,(N{-}1)/2)$.

For the normalisation constant, use Stirling's approximation
on the beta function:
\begin{equation}
  B\Bigl(\frac{1}{2},\frac{N{-}1}{2}\Bigr)
  = \frac{\sqrt{\pi}\;\Gamma\!\bigl(\frac{N-1}{2}\bigr)}
         {\Gamma\!\bigl(\frac{N}{2}\bigr)}
  \approx \sqrt{\frac{2\pi}{N-1}}
  \approx \sqrt{\frac{2\pi}{N}}\,,
\end{equation}
so $C_N \approx \sqrt{N/(2\pi)}$.

Combining:
\begin{equation}\label{eq:Psphere_full}
  P_{\rm sphere}
  \approx \sqrt{\frac{N}{2\pi}}
  \cdot \frac{2^{(N-3)/2}\,\delta^{(N-1)/2}}{(N{-}1)/2}
  = \frac{(2\delta)^{(N-1)/2}}{(N{-}1)\sqrt{2\delta}}
    \cdot \sqrt{\frac{2N}{\pi}}\,.
\end{equation}
The leading exponential is:
\begin{equation}\label{eq:lnPsphere}
  \boxed{%
  \ln P_{\rm sphere}
  = \frac{N{-}1}{2}\,\ln(2\delta)
  + O(\ln N)\,,}
\end{equation}
where $\delta = 1 - q_{\rm geom} = 4/(\beta N)$, giving
\begin{equation}\label{eq:lnPsphere_explicit}
  \ln P_{\rm sphere}
  = \frac{N{-}1}{2}\,\ln\frac{8}{\beta N}
  + O(\ln N).
\end{equation}
Since $\beta N \gg 1$, we have $\ln(8/(\beta N)) < 0$
and $P_{\rm sphere}$ is exponentially small---but the
exponent is $(N/2)\ln(8/(\beta N))$, not $-N/2$.


\subsection{Gaussian tail}

Under $q \sim \mathcal{N}(0, 1/N)$, the standardised
variable is $z = q\sqrt{N}$ and
$z_0 = q_{\rm geom}\sqrt{N}
  \approx \sqrt{N} - 4/(\beta\sqrt{N})$.
By Mill's ratio ($\bar\Phi(z) \approx \phi(z)/z$
for $z \gg 1$):
\begin{equation}\label{eq:PGauss}
  P_{\rm Gauss}
  = \bar\Phi(z_0)
  \approx \frac{e^{-z_0^2/2}}{z_0\sqrt{2\pi}}\,.
\end{equation}
The exponent:
\begin{equation}
  -\frac{z_0^2}{2}
  = -\frac{N\,q_{\rm geom}^2}{2}
  \approx -\frac{N}{2}
    \Bigl(1 - \frac{8}{\beta N}\Bigr)
  = -\frac{N}{2} + \frac{4}{\beta}\,,
\end{equation}
so
\begin{equation}\label{eq:lnPgauss}
  \boxed{%
  \ln P_{\rm Gauss}
  = -\frac{N}{2} + \frac{4}{\beta} + O(\ln N)\,.}
\end{equation}


\subsection{The discrepancy}
\label{sec:discrepancy}

Taking the ratio:
\begin{align}
  \ln\frac{P_{\rm sphere}}{P_{\rm Gauss}}
  &= \frac{N}{2}\ln\frac{8}{\beta N}
    - \Bigl(-\frac{N}{2} + \frac{4}{\beta}\Bigr)
    + O(\ln N)
  \nonumber\\[4pt]
  &= \frac{N}{2}\Bigl(\ln\frac{8}{\beta N} + 1\Bigr)
    - \frac{4}{\beta}
    + O(\ln N)
  \nonumber\\[4pt]
  &= -\frac{N}{2}\Bigl(\ln\frac{\beta N}{8} - 1\Bigr)
    + O(\ln N).
  \label{eq:ratio}
\end{align}
\begin{equation}
  \boxed{%
  \frac{P_{\rm sphere}}{P_{\rm Gauss}}
  \;\approx\;
  \exp\!\Bigl[-\frac{N}{2}\Bigl(
    \ln\frac{\beta N}{8} - 1\Bigr)\Bigr].}
\end{equation}
For $\beta N > 8e \approx 21.7$, the exponent is negative
and the sphere gives a \emph{smaller} tail than the Gaussian.
The discrepancy grows as $(N/2)\ln(\beta N)$---this is
\textbf{super-exponential} in $N$.

\paragraph{Why the Gaussian fails.}
The exact log-density is
\begin{equation}
  \ln f(q) = \text{const} + \frac{N{-}3}{2}\ln(1-q^2).
\end{equation}
Expanding around $q = 0$:
\begin{equation}\label{eq:series}
  \frac{N{-}3}{2}\ln(1-q^2)
  = -\frac{N{-}3}{2}\,q^2
  - \frac{N{-}3}{4}\,q^4
  - \frac{N{-}3}{6}\,q^6
  - \cdots
\end{equation}
The Gaussian keeps only the $q^2$ term.
For $q = O(1/\sqrt{N})$, the $q^4$ correction is $O(1)$
and can be absorbed into the normalisation.
But for $q \to 1$, \emph{every term is $O(N)$}:
$q^{2k} \approx 1$ and the $k$-th term is $-(N{-}3)/(2k)$.
The full series sums to the exact $\ln(1-q^2)$, which
diverges as $q \to 1$---qualitatively different from
the finite Gaussian exponent $-Nq^2/2 \approx -N/2$.


%% ==================================================================
\section{Geometric percolation threshold}
\label{sec:threshold}

\subsection{Expected number of geometric neighbours}

The expected number of patterns whose basins overlap with a given
pattern's basin is
\begin{equation}\label{eq:lambda}
  \lambda = (M - 1) \cdot P(q > q_{\rm geom})
  \approx e^{N\alpha} \cdot P.
\end{equation}
When $\lambda > 1$, a typical pattern has at least one
geometric neighbour; when $\lambda \gg 1$, a giant connected
component forms in the random geometric graph on the sphere.
The percolation threshold $\lambda = 1$ gives
$\alpha_c = -(\ln P)/N$.

\subsection{Gaussian prediction}

From~\eqref{eq:lnPgauss}:
\begin{equation}\label{eq:ac_gauss}
  \alpha_c^{\rm Gauss}
  = \frac{1}{2} - \frac{4}{\beta N}
  + \frac{O(\ln N)}{N}
  \;\;\xrightarrow{N \to \infty}\;\;
  \frac{1}{2}\,.
\end{equation}
This is the standard result cited for the exponential
associative memory.

\subsection{Exact spherical prediction}

From~\eqref{eq:lnPsphere_explicit}:
\begin{equation}\label{eq:ac_sphere}
  \boxed{%
  \alpha_c^{\rm sphere}
  = -\frac{1}{N}\cdot\frac{N{-}1}{2}\,\ln\frac{8}{\beta N}
  + \frac{O(\ln N)}{N}
  \;\approx\;
  \frac{1}{2}\ln\frac{\beta N}{8}
  = \frac{\ln\beta + \ln N - \ln 8}{2}\,.}
\end{equation}
This diverges as $\tfrac{1}{2}\ln N$ for any fixed $\beta > 0$.

For any fixed load $\alpha$, the condition
$\alpha > \alpha_c^{\rm sphere}$ is violated once
\begin{equation}\label{eq:Ncrit}
  N > N_{\rm crit} = \frac{8}{\beta}\,e^{2\alpha}.
\end{equation}
Beyond $N_{\rm crit}$, $\lambda \to 0$ super-exponentially
and every basin is geometrically isolated.


\subsection{Numerical comparison}

\begin{center}
\begin{tabular}{cccccc}
\toprule
$N$ & $\beta$ & $\beta N$ &
  $\alpha_c^{\rm Gauss}$ &
  $\alpha_c^{\rm sphere}$ &
  $\ln(P_{\rm sphere}/P_{\rm Gauss})$ \\
\midrule
 10 & 1 &  10 & 0.50 & 0.11 & $+1$ \\
 20 & 1 &  20 & 0.50 & 0.46 & $-1$ \\
 50 & 1 &  50 & 0.50 & 0.93 & $-22$ \\
100 & 1 & 100 & 0.50 & 1.26 & $-68$ \\
200 & 1 & 200 & 0.50 & 1.61 & $-183$ \\
\midrule
 50 & 0.2 & 10 & 0.50 & 0.11 & $+2$ \\
 50 & 0.5 & 25 & 0.50 & 0.57 & $-4$ \\
 50 & 2.0 & 100 & 0.50 & 1.26 & $-39$ \\
\bottomrule
\end{tabular}
\end{center}

For $\beta N \lesssim 20$, the two predictions are comparable
(the sphere even gives a \emph{heavier} tail than the Gaussian
for $\beta N = 10$, because the leading-order
formula~\eqref{eq:lnPsphere} drops polynomial prefactors
that matter at small~$\beta N$).
For $\beta N \gtrsim 30$, the Gaussian overestimates
$\lambda$ by many orders of magnitude.

\paragraph{Example: $\alpha = 0.5$, $\beta = 1$.}
The Gaussian predicts $\lambda \approx 1$ (percolation threshold).
The exact sphere gives $N_{\rm crit} = 8e \approx 22$.
For $N = 50$: $\lambda_{\rm sphere} \approx e^{50 \cdot 0.5
+ 25\ln(8/50)} \approx e^{25 - 46} = e^{-21} \approx 10^{-9}$.
The basin has zero geometric neighbours with overwhelming probability.


%% ==================================================================
\section{Structure of the geometric graph}
\label{sec:graph}

\subsection{When chains can form}

For $\alpha > \alpha_c^{\rm sphere}
= \tfrac{1}{2}\ln(\beta N/8)$, the geometric random
graph on $S^{N-1}$ has $\lambda > 1$ and forms a
giant connected component with chain structure:
\begin{itemize}
  \item Each hop covers angular distance
    $\sim 2\theta_c \approx 2\sqrt{2/(\beta N)}$.
  \item To traverse $O(1)$ angular distance requires
    $k \sim \sqrt{\beta N}/2$ hops.
  \item At each step, the frontier shell (annulus of width
    $\sim \theta_c$ beyond the current boundary) contains
    $\sim \lambda$ candidate patterns.
  \item After $k$ hops, the cluster contains $\sim \lambda^k$
    patterns, which is sub-exponential in~$N$ for
    $k = O(\sqrt{N})$.
\end{itemize}

\subsection{Why chains are absent at fixed $\alpha$}

For any fixed $\alpha$ and $N > N_{\rm crit}
= 8e^{2\alpha}/\beta$:
\begin{itemize}
  \item $\lambda \to 0$ super-exponentially.
  \item Each basin is an isolated island with no geometric
    neighbours.
  \item The BFS from any starting state terminates at depth~0:
    no patterns are found beyond the initial basin.
\end{itemize}

\subsection{State-space BFS at $\alpha \approx 0.5$}

Starting at $s_0 \in X_{\mu_0}$ and walking to the boundary:
\begin{itemize}
  \item $\beta N \lesssim 20$:
    the geometric graph is non-trivial; chains of a few hops
    may exist.  This is the small-$\beta N$ regime where
    the Gaussian approximation is qualitatively adequate.
  \item $\beta N \approx 22$: crossover.  $\lambda \approx 1$.
  \item $\beta N \gtrsim 30$:
    $\lambda \ll 1$; no geometric neighbours exist.
    The BFS terminates immediately.
\end{itemize}


%% ==================================================================
\section{What drives the phase transition at $\alpha = 1/2$?}
\label{sec:softmax}

The exponential associative memory
($b \to \infty$ limit of~\eqref{eq:E})
has energy $E = -\frac{1}{\beta}\ln\sum_\mu e^{\beta N\vp_\mu}$
and a well-established phase transition at $\alpha_c = 1/2$.
Since geometric percolation requires
$\alpha \gg 1/2$ for large~$N$,
this transition must be driven by a different mechanism.

\subsection{Signal-to-noise analysis}

At the target centre $\bS = \bxi^1$, the energy gradient
is proportional to a softmax-weighted sum of patterns:
\begin{equation}
  \nabla E \propto \sum_\mu w_\mu\,\bxi^\mu, \qquad
  w_\mu = \frac{e^{\beta N\vp_\mu}}{\sum_\nu e^{\beta N\vp_\nu}}.
\end{equation}
The \textbf{signal} is the target's contribution:
$e^{\beta N \cdot 1} = e^{\beta N}$.
The \textbf{noise} is the sum of all other contributions:
\begin{equation}
  \text{Noise} = \sum_{\mu \neq 1} e^{\beta N\vp_\mu}.
\end{equation}
For random overlaps $\vp_\mu \sim \mathcal{N}(0, 1/N)$
(in the Gaussian approximation), the moment generating
function gives
$\mathbb{E}[e^{\beta N\vp}] = e^{\beta^2 N/2}$.
Therefore:
\begin{equation}\label{eq:Enoise}
  \mathbb{E}[\text{Noise}]
  = (M - 1)\,e^{\beta^2 N/2}
  \approx e^{N\alpha + \beta^2 N/2}.
\end{equation}
The target dominates when
$e^{\beta N} \gg e^{N\alpha + \beta^2 N/2}$, i.e.
\begin{equation}\label{eq:retrieval_condition}
  \boxed{%
  \alpha < \beta - \frac{\beta^2}{2}
  = \beta\Bigl(1 - \frac{\beta}{2}\Bigr).}
\end{equation}
For $\beta = 1$: $\alpha_c = 1/2$.
The maximum over $\beta$ is at $\beta^* = 1$, giving
$\alpha_c^{\max} = 1/2$.

\subsection{Nature of the transition}

This is an \textbf{information-theoretic} transition:
\begin{itemize}
  \item The state does not walk to other basins.
    Each basin is geometrically isolated for $\beta N \gg 1$.
  \item Exponentially many patterns contribute to the softmax
    through their weights $e^{\beta N\vp_\mu}$, even though
    \emph{none} of them are ``active'' in the hard-wall sense
    (all have $\vp_\mu < \vp_c = 1 - 1/(\beta N)$ at the
    target centre).
  \item When $\alpha > 1/2$, the collective noise from
    $\sim e^{N/2}$ patterns overwhelms the target signal.
    The target ceases to be a stable fixed point of the dynamics.
\end{itemize}

\paragraph{Key distinction from $b = O(1)$.}
For $b = O(1)$ (the LSR regime), the hard wall at $\vp_c = O(1)$
creates genuine infinite barriers between well-separated basins.
The phase transition occurs when the state can percolate
through a network of corridors.
For $b = \beta N$, the hard wall at $\vp_c = 1 - O(1/N)$ is
so restrictive that basins are microscopic and isolated.
The transition occurs because the softmax energy is effectively
a global interaction---every pattern contributes everywhere on
the sphere through its exponential weight, regardless of the
hard-wall constraint.

\subsection{Spherical correction to the signal-to-noise threshold}

The Gaussian MGF $\mathbb{E}[e^{\beta N q}] = e^{\beta^2 N/2}$
receives corrections from the exact spherical distribution.
The exact MGF is
\begin{equation}\label{eq:MGF_sphere}
  \mathbb{E}_{\rm sphere}[e^{\beta N q}]
  = C_N \int_{-1}^{1} e^{\beta N q}\,(1-q^2)^{(N-3)/2}\,dq.
\end{equation}
For large $N$, evaluate by saddle point.  The exponent is
$g(q) = \beta N q + \tfrac{N-3}{2}\ln(1 - q^2)$,
with $g'(q) = \beta N - (N{-}3)q/(1-q^2) = 0$, giving
\begin{equation}\label{eq:saddle}
  q^* = \frac{\beta(1 - q^{*2})}{1 - 3/N}
  \;\approx\; \beta(1 - q^{*2}).
\end{equation}
For $\beta < 1$: $q^* \in (0, 1)$ and the saddle is in the
interior.
For $\beta = 1$: $q^* = (\sqrt{5}-1)/2 \approx 0.618$
(the golden ratio minus one).
The saddle-point exponent $g(q^*)$ determines the leading
behaviour of the MGF.  When $g(q^*)$ differs from
$\beta^2 N/2$ (the Gaussian result), the signal-to-noise
threshold $\alpha_c$ shifts.
A systematic evaluation of this correction is an important
direction for future work.


%% ==================================================================
\section{Finite-$\beta N$ crossover}
\label{sec:crossover}

The results of the preceding sections delineate two regimes:

\begin{center}
\begin{tabular}{lcc}
\toprule
  & Small $\beta N$ & Large $\beta N$ \\
\midrule
Basins & wide caps, $\theta_c = O(1)$
       & microscopic, $\theta_c \sim 1/\sqrt{\beta N}$ \\
Geometric overlaps & abundant
                   & exponentially rare \\
Energy barriers & $O(N/b)$, extensive
               & $O(1/\beta)$, thermal \\
Connectivity bottleneck & energy
                        & geometry \\
Phase transition mechanism & basin percolation
                          & softmax signal-to-noise \\
$\alpha_c$ & $\frac{1}{2}(1-1/b)^2$
           & $1/2$ (S/N), or $\frac{1}{2}\ln(\beta N/8)$ (geom.) \\
\bottomrule
\end{tabular}
\end{center}

The crossover between the two regimes is determined by
the condition
\begin{equation}
  \alpha_c^{\rm sphere}(\beta, N) = \alpha,
  \qquad\text{i.e.}\qquad
  \beta N = 8\,e^{2\alpha}.
\end{equation}
At $\alpha = 0.5$: $\beta N_{\rm cross} \approx 22$.
At $\alpha = 1.0$: $\beta N_{\rm cross} \approx 59$.
At $\alpha = 2.0$: $\beta N_{\rm cross} \approx 436$.

\paragraph{Finite-$N$ corrections.}
The leading formula~\eqref{eq:ac_sphere} drops polynomial
prefactors.  Retaining them from~\eqref{eq:Psphere_full}:
\begin{equation}\label{eq:ac_refined}
  \alpha_c^{\rm sphere}
  = \frac{1}{2}\ln\frac{\beta N}{8}
  + \frac{1}{N}\ln\sqrt{\frac{\pi\,\beta N}{4}}
  + O\!\Bigl(\frac{\ln N}{N}\Bigr).
\end{equation}
The subleading correction is positive and $O(\ln N / N)$,
slightly increasing $\alpha_c$ at finite~$N$.


%% ==================================================================
\section{Summary}
\label{sec:summary}

\begin{enumerate}
  \item \textbf{Geometry first, energy second.}
    For $b = \beta N$, energy barriers between geometrically
    overlapping basins are $O(1/\beta)$---thermally irrelevant.
    The question reduces to: do any caps overlap?

  \item \textbf{The Gaussian approximation fails catastrophically
    for $\beta N \gg 1$.}
    At $q_{\rm geom} = 1 - 4/(\beta N)$,
    the Gaussian overestimates the tail probability by
    $\exp\!\bigl[\tfrac{N}{2}(\ln(\beta N/8) - 1)\bigr]$,
    a super-exponential factor.

  \item \textbf{Geometric percolation requires
    $\alpha \sim \tfrac{1}{2}\ln N$.}
    For any fixed load $\alpha$, the geometric graph has
    no edges once $N > 8e^{2\alpha}/\beta$.
    At $\alpha = 0.5$, $\beta = 1$: no neighbours for $N > 22$.

  \item \textbf{The transition at $\alpha = 1/2$ is not
    geometric percolation.}
    It is driven by signal-to-noise competition in the softmax:
    $e^{N/2}$ individually negligible patterns collectively
    overwhelm the target.

  \item \textbf{Crossover at $\beta N \approx 8\,e^{2\alpha}$.}
    Below this, the percolation picture from the $b = O(1)$
    analysis applies qualitatively.
    Above it, basins are isolated and the softmax mechanism
    dominates.
\end{enumerate}


\end{document}
